\section{Evaluación y evidencia}

\subsection{Benchmark Stack}

El objetivo de la optimización de preferencias es que el modelo muestre una mejora consistente en diversos benchmarks. Los niveles de evaluación habituales son:

\begin{table}[H]
\centering
\begin{tabular}{p{0.2\textwidth} p{0.28\textwidth} p{0.34\textwidth}}
\toprule
Nivel & Métricas de ejemplo & Descripción \\
\midrule
Benchmark automatizado & HumanEval, MMLU, TruthfulQA & Cuantifica las capacidades y el reasoning \\
Consistencia de preferencias & Pairwise Preference Accuracy & Si el modelo de recompensa puede predecir las preferencias humanas \\
Verificación de alineación & Red Team, Safety Gym & Desempeño del modelo dentro de los límites de seguridad \\
Métricas de despliegue & Rejection Rate, Hallucination Rate & Calidad de la experiencia en el uso real \\
\bottomrule
\end{tabular}
\caption{Pila de evaluación de RLHF/DPO}
\end{table}

\subsection{Evidence Matrix}

Antes del despliegue hay que organizar una evidence matrix para construir la cadena de evidencia de la alineación y la validación:

\begin{table}[H]
\centering
\begin{tabular}{p{0.2\textwidth} p{0.24\textwidth} p{0.24\textwidth} p{0.24\textwidth}}
\toprule
Evidence & Source & Metric & Decision Use \\
\midrule
Pairwise human ranking & crowdsourced label studio & Accuracy & Reward model calibration \\
Automated guardrails & prompt test suite & Reject rate & Safety filter tuning \\
Long-context consistency & multi-turn benchmark & Drift score & Deployment guardrail \\
Deployment logs & tooling + telemetry & Incident count & Postmortem loop \\
\bottomrule
\end{tabular}
\caption{Evidence matrix para la validación entre etapas}
\end{table}

\subsection{Calibration y Robustness}

No basta con alinear la evaluation: también hay que mantener la estabilidad ante distintas inyecciones semánticas, longitudes de contexto y variaciones del token budget. Las medidas habituales son:

\begin{itemize}
    \item Aplicar temperature calibration a la salida del reward model para llevarla al intervalo 0-1
    \item Construir un multi-context benchmark para detectar el evidence drift
    \item Usar counterfactual prompts (e.g. `change numeric context`) para verificar que la tendencia de cambio del reward sea consistente
    \item Registrar los rejection case IDs para poder hacer drill-down
\end{itemize}

\begin{warningbox}{Riesgo de juzgar mal el drift del modelo de preferencias}
Aunque el reward model tenga buen desempeño en el conjunto de entrenamiento, puede presentar drift con prompts reales. Sin calibration ni drift detection, no puede garantizarse que la alignment siga vigente después del despliegue.
\end{warningbox}

\subsection{Métricas de despliegue y evidence ranking}

Cuando el modelo entra en despliegue gradual o completo, la evaluation stack debe conectarse de forma robusta con la capa de monitoreo:
\begin{itemize}
    \item \textbf{Rejection/deferral rate}: determina si la guardrail es demasiado permisiva o demasiado estricta.
    \item \textbf{Hallucination trend}: agrupar por topic y revisar si la hallucination rate se concentra en direcciones específicas.
    \item \textbf{Latency drift}: si, después de actualizar la política, la latency rebasa el SLO.
    \item \textbf{Evidence score}: ordenar por niveles el benchmark, la human eval y los deployment logs, para dar al decision board un fundamento en tiempo real.
\end{itemize}

\begin{knowledgebox}{El valor del evidence ranking}
Asignar pesos distintos a la evidence de cada etapa (por ejemplo: logs de despliegue > benchmark) permite avanzar con rapidez en la iteración de versiones sin sacrificar la seguridad.
\end{knowledgebox}

\begin{importantbox}{El tablero de evaluación necesita visualización en tiempo real}
Archit advierte: \textit{«La evidence matrix no es una tabla escrita en un documento, sino un mapa de calor que se actualiza continuamente en el dashboard»}; solo la visualización en tiempo real permite descubrir los puntos ciegos del drift.
\end{importantbox}

\subsection{Resumen de la sección}

El sistema de evaluación y evidencia debe cubrir, nivel por nivel, desde el benchmark hasta el despliegue; la Evidence Matrix es el eje que combina distintas señales en conclusiones accionables, y las deploy metrics extienden la evaluation stack hasta production.
